% ============================================================================
% Algoritmos 4, 5 y 6 -- formalismo PARDO unificado (feedback entrega parcial)
% Preámbulo requerido (mismo estilo del informe parcial):
%   \usepackage[ruled,linesnumbered,noend]{algorithm2e}
%   \SetKwFor{ForAllPar}{forall}{pardo}{end}
%   \SetKwInput{KwLocal}{Datos locales}
%   \SetKwInput{KwShared}{Memoria global}
%   \SetKwComment{tcp}{$\triangleright$\ }{}
%   \SetKwFunction{Bcast}{Bcast}\SetKwFunction{Scatterv}{Scatterv}
%   \SetKwFunction{Gather}{Gather}\SetKwFunction{Gatherv}{Gatherv}
%   \SetKwFunction{Predict}{knn\_predict}
% Convención única: TODO bucle concurrente se escribe
%   forall <índice> <- <ini> to <fin> pardo
% tanto para procesadores PRAM (Alg. 4), procesos MPI (Algs. 5 y 6) como para
% consultas locales independientes (paralelismo de datos dentro del proceso).
% ============================================================================

% ---------------------------------------------------------------- Algoritmo 4
\begin{algorithm}[H]
\caption{KNN en PRAM CREW}\label{alg:knn-pram}
\KwIn{Datos $X\in\mathbb{R}^{(n_{\text{train}}+m_{\text{test}})\times d}$ y etiquetas
      $y\in\{0,\dots,9\}^{n_{\text{train}}+m_{\text{test}}}$}
\KwOut{Exactitud del modelo y tiempo de ejecución}
\KwShared{$X_{\text{tr}}, y_{\text{tr}}$ (CR); $\hat{y}$ (EW)}
\BlankLine
$X_{\text{tr}}, X_{\text{te}}, y_{\text{tr}}, y_{\text{te}} \leftarrow
  \texttt{train\_test\_split}(X, y, \texttt{test\_size}=0.2, \texttt{random\_state}=42)$\;
$k \leftarrow 3$\;
$t_{\text{inicio}} \leftarrow \texttt{time}()$\;
$\hat{y} \leftarrow 0 \in \{0,\dots,9\}^{m_{\text{test}}}$\;
\ForAllPar{$j \leftarrow 1$ \KwTo $m_{\text{test}}$}{
  $\hat{y}[j] \leftarrow$ \Predict{$X_{\text{te}}[j], X_{\text{tr}}, y_{\text{tr}}, k$}
    \tcp*{procesador $j$; Algoritmo 2}
}
$\text{exactitud} \leftarrow \texttt{mean}(\hat{y} = y_{\text{te}})$\;
$t_{\text{fin}} \leftarrow \texttt{time}()$\;
\Return{exactitud, $t_{\text{fin}} - t_{\text{inicio}}$}\;
\end{algorithm}

% ---------------------------------------------------------------- Algoritmo 5
\begin{algorithm}[H]
\caption{KNN distribuido con MPI: replicación de $X_{\text{tr}}$, partición de $X_{\text{te}}$}\label{alg:knn-mpi}
\KwIn{Datos $X$, etiquetas $y$, $s$ procesos, $k=3$}
\KwOut{Exactitud del modelo y tiempo de ejecución (en el proceso $0$)}
\KwLocal{El proceso $p$ recibe $m_p=\lfloor m_{\text{test}}/s\rfloor+[\,p<m_{\text{test}}\bmod s\,]$
         consultas desde $\mathrm{displ}_p=\sum_{i<p} m_i$}
\BlankLine
\ForAllPar{$p \leftarrow 0$ \KwTo $s-1$}{
  \If{$p = 0$}{
    $X_{\text{tr}}, X_{\text{te}}, y_{\text{tr}}, y_{\text{te}} \leftarrow
      \texttt{train\_test\_split}(X, y, \texttt{test\_size}=0.2, \texttt{random\_state}=42)$\;
    $t_{\text{inicio}} \leftarrow \texttt{Wtime}()$\;
  }
  \tcp*[l]{replicación: $\log s\,(\alpha + n_{\text{train}} d\,\beta)$}
  $X_{\text{tr}}, y_{\text{tr}} \leftarrow$ \Bcast{$X_{\text{tr}}, y_{\text{tr}}$, raíz $=0$}\;
  \tcp*[l]{partición de consultas: $\alpha\log s + m_{\text{test}} d\,\beta$}
  $X^{(p)}_{\text{te}} \leftarrow$ \Scatterv{$X_{\text{te}}$, $(m_i d)_i$, $(\mathrm{displ}_i d)_i$, raíz $=0$}\;
  \tcp*[l]{cómputo local: $m_p\, n_{\text{train}}\,(3d)$ FLOPs}
  \ForAllPar{$j \leftarrow 1$ \KwTo $m_p$}{
    $\hat{y}^{(p)}[j] \leftarrow$ \Predict{$X^{(p)}_{\text{te}}[j], X_{\text{tr}}, y_{\text{tr}}, k$}
      \tcp*{Algoritmo 2; CR}
  }
  \tcp*[l]{recolección: $\alpha\log s + m_{\text{test}}\,\beta$}
  $\hat{y}[\mathrm{displ}_p+1 \,..\, \mathrm{displ}_p+m_p] \leftarrow$
     \Gatherv{$\hat{y}^{(p)}$, $(m_i)_i$, $(\mathrm{displ}_i)_i$, raíz $=0$}
     \tcp*{EW: bloques disjuntos}
  \If{$p = 0$}{
    $\text{exactitud} \leftarrow \texttt{mean}(\hat{y} = y_{\text{te}})$\;
    $t_{\text{fin}} \leftarrow \texttt{Wtime}()$\;
    \Return{exactitud, $t_{\text{fin}}-t_{\text{inicio}}$}\;
  }
}
\end{algorithm}

\paragraph{Formalismo \texttt{pardo}.}
En el Algoritmo~\ref{alg:knn-pram} el índice $j$ recorre los $\mte$ procesadores de una PRAM
CREW~\cite{jaja1992}: todos leen concurrentemente $\Xtr$ y $\ytr$ (CR) y el procesador $j$ escribe
sólo en $\hat{y}[j]$ (EW). En los Algoritmos~\ref{alg:knn-mpi} y~\ref{alg:knn-dag}, el bucle
externo \textbf{forall} $p\leftarrow 0$ \textbf{to} $s-1$ \textbf{pardo} expresa el paradigma
SPMD: los $s$ procesos ejecutan concurrentemente el mismo cuerpo y se distinguen sólo por su rango;
las colectivas son puntos de sincronización implícita. El bucle interno expresa que las consultas
locales son independientes (no hay aristas entre ellas en el DAG). Como cada proceso ocupa un único
núcleo, ese \texttt{pardo} interno se ejecuta como paralelismo de datos ---producto matricial por
bloques--- y, por el teorema de Brent~\cite{brent1974},
\begin{equation}
  \Tcomp(s)\le\frac{W}{s}+T_\infty .\label{eq:brent}
\end{equation}
Con $m_p$ procesadores virtuales por proceso ($s\,m_p=\mte$), el Algoritmo~\ref{alg:knn-mpi} se
reduce exactamente al Algoritmo~\ref{alg:knn-pram}.


% ---------------------------------------------------------------- Algoritmo 6
\begin{algorithm}[H]
\caption{KNN distribuido con MPI según el DAG: partición de $X_{\text{tr}}$, difusión de $X_{\text{te}}$}\label{alg:knn-dag}
\KwIn{Datos $X$, etiquetas $y$, $s$ procesos, $k=3$}
\KwOut{Exactitud del modelo y tiempo de ejecución (en el proceso $0$)}
\KwLocal{El proceso $p$ recibe $n_p=\lfloor n_{\text{train}}/s\rfloor+[\,p<n_{\text{train}}\bmod s\,]$
         puntos de entrenamiento $(X^{(p)}_{\text{tr}}, y^{(p)}_{\text{tr}})$}
\BlankLine
\ForAllPar{$p \leftarrow 0$ \KwTo $s-1$}{
  \If{$p = 0$}{
    $X_{\text{tr}}, X_{\text{te}}, y_{\text{tr}}, y_{\text{te}} \leftarrow
      \texttt{train\_test\_split}(X, y, \texttt{test\_size}=0.2, \texttt{random\_state}=42)$\;
    $t_{\text{inicio}} \leftarrow \texttt{Wtime}()$\;
  }
  \tcp*[l]{partición del entrenamiento: $\alpha\log s + n_{\text{train}} d\,\beta$}
  $X^{(p)}_{\text{tr}}, y^{(p)}_{\text{tr}} \leftarrow$ \Scatterv{$X_{\text{tr}}, y_{\text{tr}}$, raíz $=0$}\;
  \tcp*[l]{difusión de consultas: $\log s\,(\alpha + m_{\text{test}} d\,\beta)$}
  $X_{\text{te}} \leftarrow$ \Bcast{$X_{\text{te}}$, raíz $=0$}\;
  \tcp*[l]{cómputo local: $m_{\text{test}}\, n_p\,(3d)$ FLOPs}
  \ForAllPar{$j \leftarrow 1$ \KwTo $m_{\text{test}}$}{
    $(\delta^{(p)}_j, \ell^{(p)}_j) \leftarrow$ $k$ menores distancias de $X_{\text{te}}[j]$ a $X^{(p)}_{\text{tr}}$ y sus etiquetas\;
  }
  \tcp*[l]{recolección de candidatos: $\alpha\log s + m_{\text{test}}\,k\,s\,\beta$}
  $\{(\delta^{(i)}, \ell^{(i)})\}_{i=0}^{s-1} \leftarrow$ \Gather{$(\delta^{(p)}, \ell^{(p)})$, raíz $=0$}\;
  \If{$p = 0$}{
    \For{$j \leftarrow 1$ \KwTo $m_{\text{test}}$}{
      $\hat{y}[j] \leftarrow \texttt{moda}\big(k$ etiquetas de menor distancia entre los $s\,k$ candidatos de $j\big)$
        \tcp*{fusión serial $O(s\,k)$}
    }
    $\text{exactitud} \leftarrow \texttt{mean}(\hat{y} = y_{\text{te}})$\;
    $t_{\text{fin}} \leftarrow \texttt{Wtime}()$\;
    \Return{exactitud, $t_{\text{fin}}-t_{\text{inicio}}$}\;
  }
}
\end{algorithm}

% ---------------------------------------------------------------- Texto sugerido
% \paragraph{Formalismo \texttt{pardo}.}
% Los tres algoritmos paralelos usan la misma construcción
% \textbf{forall} $i\leftarrow a$ \textbf{to} $b$ \textbf{pardo} [JáJá]. En el
% Algoritmo~\ref{alg:knn-pram} el índice recorre los $m_{\text{test}}$
% procesadores de la PRAM CREW. En los Algoritmos~\ref{alg:knn-mpi} y
% \ref{alg:knn-dag} el bucle externo sobre $p$ expresa el paradigma SPMD: los
% $s$ procesos MPI ejecutan concurrentemente el mismo cuerpo y se distinguen sólo
% por su rango; las colectivas son puntos de sincronización implícita. El bucle
% interno expresa que las consultas son independientes (no existen aristas entre
% ellas en el DAG); como cada proceso ocupa un único núcleo, se ejecuta como
% paralelismo de datos (producto matricial por bloques, BLAS/SIMD) y, por el
% teorema de Brent, $T_{\text{comp}}(s)\le W/s+T_\infty$. Con $m_p$ procesadores
% virtuales por proceso ($s\,m_p=m_{\text{test}}$), el Algoritmo~\ref{alg:knn-mpi}
% se reduce exactamente al Algoritmo~\ref{alg:knn-pram}.
%
% \paragraph{Corrección del Algoritmo~\ref{alg:knn-dag}.}
% Si $x$ pertenece a los $k$ vecinos globales de $q$ y $x\in X^{(p)}_{\text{tr}}$,
% entonces $x$ también está entre los $k$ vecinos de $q$ en el bloque $p$; luego
% la unión de los $s\,k$ candidatos contiene a los $k$ vecinos globales y la
% predicción coincide con la secuencial (salvo empates exactos de distancia).
%
% \paragraph{Comparación de costos.}
% Alg.~5: $T_{\text{comm}}=O\big(\alpha\log s+\beta(n_{\text{train}}d\log s+m_{\text{test}}d)\big)$.
% Alg.~6: $T_{\text{comm}}=O\big(\alpha\log s+\beta(n_{\text{train}}d+m_{\text{test}}d\log s+m_{\text{test}}ks)\big)$,
% más una fase serial $O(m_{\text{test}}\,s\,k)$ en el proceso $0$.
% Como $n_{\text{train}}=4\,m_{\text{test}}$, el Alg.~6 transfiere menos datos
% en la difusión, pero su recolección y su fusión crecen linealmente con $s$.
